\documentclass[11pt,a4paper,oneside]{report}
\usepackage[spanish,activeacute,es-noshorthands]{babel}
\usepackage[utf8]{inputenc}
\usepackage[T1]{fontenc}
\usepackage{lmodern}
\usepackage{amsmath, amssymb}
\usepackage{graphicx}
\usepackage[dvipsnames]{xcolor}
\usepackage{geometry}
\geometry{margin=2.3cm, headheight=14pt}
\usepackage{microtype}
\usepackage{booktabs}
\usepackage{longtable}
\usepackage{enumitem}
\usepackage{fancyhdr}
\usepackage{tcolorbox}
\tcbuselibrary{skins, breakable}
\usepackage{caption}
\captionsetup{font=small, labelfont=bf}
\usepackage{titlesec}
\usepackage{hyperref}
\hypersetup{
  colorlinks=true, linkcolor=NavyBlue, urlcolor=BrickRed,
  pdftitle={Plan de documentacion CubeSat},
  pdfauthor={INTISAT --- CubeSat-EdgeAI-Payload}
}

\definecolor{intisat}{RGB}{26, 54, 93}
\definecolor{intisataccent}{RGB}{197, 143, 0}
\definecolor{ok}{RGB}{47, 133, 90}
\definecolor{warn}{RGB}{197, 60, 60}
\definecolor{prio}{RGB}{197, 143, 0}

\titleformat{\chapter}[hang]
  {\normalfont\Huge\bfseries\color{intisat}}
  {\thechapter.}{0.6em}{}
\titlespacing*{\chapter}{0pt}{-20pt}{20pt}

\pagestyle{fancy}
\fancyhf{}
\fancyhead[L]{\nouppercase{\leftmark}}
\fancyhead[R]{\thepage}
\renewcommand{\headrulewidth}{0.4pt}

\newtcolorbox{tip}[1][]{
  colback=green!5, colframe=ok,
  fonttitle=\bfseries, title=Recomendacion,
  breakable, sharp corners=south, #1
}
\newtcolorbox{nota}[1][]{
  colback=blue!4, colframe=intisat,
  fonttitle=\bfseries, title=Nota,
  breakable, sharp corners=south, #1
}
\newtcolorbox{prioritario}[1][]{
  colback=orange!5, colframe=intisataccent,
  fonttitle=\bfseries, title=Manual prioritario,
  breakable, sharp corners=south, #1
}

\begin{document}

\begin{titlepage}
\centering
\vspace*{1cm}
{\color{intisat}\rule{\textwidth}{2pt}}
\vspace{0.6cm}

{\Huge\bfseries\color{intisat} Plan de documentacion}\\[0.4em]
{\LARGE\bfseries Coleccion de manuales para construir un CubeSat}

\vspace{0.7cm}
{\Large\itshape Dimensionamiento, contenidos y prioridades\\
para una serie tecnica completa}

\vspace{1cm}
{\color{intisataccent}\rule{0.6\textwidth}{1pt}}

\vspace{2cm}

\begin{tabular}{rl}
\textbf{Proyecto:} & CubeSat-EdgeAI-Payload \\
\textbf{Plataforma:} & INTISAT \\
\textbf{Audiencia:} & Equipo tecnico, estudiantes, futuras misiones \\
\textbf{Total proyectado:} & 11 manuales --- aprox. 510 paginas \\
\textbf{Repositorio:} & \texttt{github.com/JaminYC/CubeSat-EdgeAI-Payload} \\
\textbf{Fecha:} & \today \\
\textbf{Documento:} & v1.0 --- plan inicial
\end{tabular}

\vfill

\begin{minipage}{0.85\textwidth}
\textbf{Resumen.}\quad
Plan estructurado para producir una coleccion completa de manuales
tecnicos que cubre el ciclo de vida de un CubeSat: desde los fundamentos
fisicos y de subsistemas hasta integracion, testing y operaciones. El
plan dimensiona 11 manuales (7 nuevos y 4 ya existentes) con sus
contenidos, prioridades, dependencias de lectura y tiempos estimados de
escritura. Permite al equipo decidir cuales redactar primero segun el
estado del proyecto.
\end{minipage}

\vspace{1.2cm}
{\color{intisat}\rule{\textwidth}{2pt}}
\end{titlepage}

\tableofcontents
\clearpage

% =====================================================================
\chapter{Como esta organizado este plan}

Un CubeSat es un sistema complejo. Aprenderlo requiere documentos que
combinen \emph{fundamentos teoricos} con \emph{aplicaciones practicas} y
\emph{ejemplos del INTISAT}.

Esta coleccion se organiza en 4 niveles:

\begin{itemize}[leftmargin=2em]
  \item \textbf{Nivel 0 --- Introduccion}: que es un CubeSat, ecosistema,
        actores, fases.
  \item \textbf{Nivel 1 --- Subsistemas}: un manual por subsistema (potencia,
        comms, ADCS, etc.).
  \item \textbf{Nivel 2 --- Cross-cutting}: temas que atraviesan todos los
        subsistemas (testing, operaciones, system engineering).
  \item \textbf{Nivel 3 --- Especifico de mision}: lo unico del INTISAT
        (payload de microscopia, OLED, FPM).
\end{itemize}

Cada manual sigue la \textbf{misma plantilla} (la misma de
\texttt{Manual\_Electronica\_Practica.pdf} y \texttt{Curso\_Vision\_Computadora.pdf}):

\begin{enumerate}[leftmargin=2em]
  \item Prefacio con audiencia y prerequisitos.
  \item Cuerpo en capitulos con: \emph{por que importa, intuicion, matematica,
        conexion con el payload, cuidado, ejercicios}.
  \item Apendices con tablas de consulta y referencias.
  \item Bibliografia.
\end{enumerate}

\section{Lo que ya existe}

\begin{center}
\begin{tabular}{p{4.5cm}p{1.5cm}p{2cm}p{5cm}}
\toprule
\textbf{Documento} & \textbf{Pag} & \textbf{Tipo} & \textbf{Estado} \\
\midrule
Curso\_Vision\_Computadora.pdf      & 81  & libro    & completo \\
Manual\_Electronica\_Practica.pdf   & 44  & manual   & completo \\
Raspberry\_Pi\_5\_Profundo.pdf      & 25  & manual   & completo \\
Pipeline\_IA\_Microscopia.pdf       & 13  & paper    & completo \\
Informe\_Mascaras\_Lensless.pdf     & 21  & informe  & completo \\
Bibliografia\_AI\_Espacial.pdf      & 18  & referencia & completo \\
Referencias\_OpenSource\_CubeSat.pdf & 18  & referencia & completo \\
\bottomrule
\end{tabular}
\end{center}

\textbf{Total existente: 220 paginas}.

\section{Lo que falta (este plan propone)}

\begin{center}
\begin{tabular}{lcccc}
\toprule
\textbf{Manual} & \textbf{Pag} & \textbf{Prio} & \textbf{Tiempo} & \textbf{Audiencia} \\
\midrule
M0: CubeSat 101                      & 50 & bigstarbigstarbigstar & 1 sem & todo el equipo \\
M1: Mecanica y estructura            & 50 & bigstarbigstar  & 1 sem & ing mecanica \\
M2: Sistemas de potencia (EPS)       & 60 & bigstarbigstarbigstar & 1.5 sem & elec/potencia \\
M3: Control de actitud (ADCS)        & 60 & bigstarbigstar  & 2 sem & control \\
M4: Comunicaciones (RF + antenas)    & 60 & bigstarbigstarbigstar & 1.5 sem & RF \\
M5: Flight software y OBC            & 60 & bigstarbigstarbigstar & 1.5 sem & software \\
M6: Control termico                  & 40 & bigstarbigstar  & 1 sem & termico \\
M7: Integracion y test campaign      & 60 & bigstarbigstarbigstar & 1.5 sem & AIT \\
M8: Operaciones y ground station     & 50 & bigstarbigstar  & 1 sem & ops \\
M9: Project management y SE          & 40 & bigstar   & 1 sem & lider proyecto \\
M10: CM5 + carrier para INTISAT      & 50 & bigstarbigstarbigstar & 1.5 sem & hw payload \\
\bottomrule
\end{tabular}
\end{center}

\textbf{Total proyectado nuevo: 580 paginas} --- ~15 semanas de escritura
(repartido entre el equipo: 4 personas $\to$ 4 semanas reales).

% =====================================================================
\chapter{Manual M0: CubeSat 101 --- introduccion al ecosistema}

\begin{prioritario}
Prioridad: $\bigstar\bigstar\bigstar$ (top). Es el manual que abre la
serie y todo el equipo deberia leer primero.
\end{prioritario}

\textbf{Paginas estimadas}: 50

\textbf{Audiencia}: cualquier persona nueva en el proyecto, sin background
aeroespacial.

\textbf{Objetivo}: que despues de leerlo, la persona pueda mantener una
conversacion tecnica con un ingeniero de CubeSat y entender los grandes
bloques.

\section{Contenido propuesto (TOC)}

\begin{enumerate}[leftmargin=2em]
  \item \textbf{Que es un CubeSat}: definicion CalPoly, formatos (1U, 3U, 6U, 12U),
        estandar de despliegue.
  \item \textbf{Historia breve}: 1999 origen, GeneSat-1, PhoneSat, PhiSat-1,
        constelaciones modernas.
  \item \textbf{Anatomia de un CubeSat}: los 6 subsistemas (mecanica, EPS,
        ADCS, COMMS, OBC, payload).
  \item \textbf{Fases de vida del satelite}: concepto, diseño, integracion,
        lanzamiento, operacion, decomission.
  \item \textbf{Como llega al espacio}: rideshare, lanzadores tipicos (SpaceX
        Transporter, Rocket Lab, ISRO PSLV).
  \item \textbf{Despliegue desde el lanzador}: P-POD, NanoRacks, ISS deployer.
  \item \textbf{Orbita LEO 101}: altura, periodo, eclipses, regimen termico,
        radiacion.
  \item \textbf{El ecosistema}: NASA, ESA, agencias espaciales LATAM,
        universidades, startups NewSpace.
  \item \textbf{Estandares clave}: CalPoly CDS, ECSS, NASA GEVS,
        MIL-STD-810.
  \item \textbf{Vocabulario}: glosario de 100+ siglas (TRL, FDIR, AIT, RF,
        BOL/EOL, etc.).
\end{enumerate}

\textbf{Por que es prioritario}: sin esto, los manuales subsiguientes no
tienen contexto. Es como un \textbf{onboarding} para un miembro nuevo.

% =====================================================================
\chapter{Manual M1: Mecanica y estructura}

\begin{prioritario}
Prioridad: $\bigstar\bigstar$ (medio).
\end{prioritario}

\textbf{Paginas}: 50 ~|~ \textbf{Audiencia}: ingenieria mecanica.

\section{Contenido propuesto}

\begin{enumerate}[leftmargin=2em]
  \item \textbf{Estandares CalPoly}: dimensiones, masa, centro de gravedad,
        rails de aluminio anodizado.
  \item \textbf{Materiales}: aluminio 6061-T6, titanio, magnesio,
        composites de fibra de carbono.
  \item \textbf{Diseño estructural}: cargas estaticas, dinamicas, modos
        propios, factor de seguridad.
  \item \textbf{Analisis de elementos finitos (FEA)}: que es, herramientas
        (ANSYS, SimScale, FreeCAD-FEM).
  \item \textbf{Tornilleria espacial}: M2, M3 con loctite, retenes, locking.
  \item \textbf{Despliegues}: antenas, paneles solares, mecanismos burn-wire.
  \item \textbf{PC-104}: el estandar mecanico/electrico, dimensiones, pinout.
  \item \textbf{Test de ajuste y CG}: como verificar el centro de masa antes
        del lanzamiento.
  \item \textbf{Outgassing}: por que importa, materiales prohibidos
        (UV degradable, plasticos con plastificantes).
\end{enumerate}

% =====================================================================
\chapter{Manual M2: Sistemas de potencia (EPS)}

\begin{prioritario}
Prioridad: $\bigstar\bigstar\bigstar$ (top --- la energia es restriccion
critica de todo CubeSat).
\end{prioritario}

\textbf{Paginas}: 60 ~|~ \textbf{Audiencia}: electronica de potencia.

\section{Contenido propuesto}

\begin{enumerate}[leftmargin=2em]
  \item \textbf{Generacion solar}: celdas GaAs vs Si, eficiencia (28\%
        triple-juncion vs 15\% Si), degradacion en orbita.
  \item \textbf{Calculo de generacion}: solar constant, angulo de incidencia,
        sombras, eclipse fraction.
  \item \textbf{Maximum Power Point Tracking (MPPT)}: por que se necesita,
        algoritmos (P\&O, IncCond).
  \item \textbf{Bateria Li-ion}: quimica, voltajes, capacidad, degradacion
        ciclica, depth of discharge (DoD).
  \item \textbf{Battery Management System (BMS)}: balance, proteccion contra
        sobre/sub-voltaje, sobre-corriente, temperatura.
  \item \textbf{Reguladores}: LDO vs switching, eficiencia, ruido,
        ripple.
  \item \textbf{Distribucion}: rails (3.3V, 5V, 12V), proteccion contra
        cortos (PTC, eFuse), enable lines.
  \item \textbf{Power budget}: como armar la tabla, modos de mision,
        margenes ($\geq 20\%$).
  \item \textbf{EPS comerciales}: GomSpace P31u, ISIS iEPS, EnduroSat
        EPS-1, especificaciones y precios.
  \item \textbf{Tests de potencia}: validacion en banco, ciclos de carga,
        TVAC.
\end{enumerate}

\section{Conexion al INTISAT}

Tu payload ya tiene un presupuesto teorico (\texttt{Energy\_Budget\_Microscopia.xlsx}).
Este manual lo lleva al \textbf{nivel mision}: como dimensionar paneles +
bateria + reguladores para alimentar tu payload + todos los otros
subsistemas durante una vida util de 1 año en orbita LEO.

% =====================================================================
\chapter{Manual M3: Control de actitud (ADCS)}

\textbf{Paginas}: 60 ~|~ \textbf{Audiencia}: control / aeroespacial.
~|~ Prioridad: $\bigstar\bigstar$.

\section{Contenido}

\begin{enumerate}[leftmargin=2em]
  \item \textbf{Por que se necesita}: apuntar paneles al sol, antena a la
        tierra, payload al objetivo.
  \item \textbf{Sensores de actitud}: gyros, magnetometros, sun sensors,
        star trackers, GPS.
  \item \textbf{Actuadores}: magnetorquers, reaction wheels, thrusters
        (raros en CubeSat).
  \item \textbf{Algoritmos}: B-dot (detumble), TRIAD, Kalman filter, PID,
        LQR.
  \item \textbf{Determinacion de actitud}: usar 2 sensores no paralelos
        (sol + magnetometro tipico).
  \item \textbf{Control de actitud}: closed loop, deadbands, slew rates.
  \item \textbf{Modos}: detumble, sun-pointing, nadir, target, idle.
  \item \textbf{Perturbaciones}: gradiente gravitacional, presion solar,
        drag atmosferico.
  \item \textbf{Componentes comerciales}: CubeWheel, NSS, MTQ, Cubespace,
        BlueCanyon XACT.
  \item \textbf{Validacion en lab}: air bearing, Helmholtz cage, sol
        simulado.
\end{enumerate}

% =====================================================================
\chapter{Manual M4: Comunicaciones}

\begin{prioritario}
Prioridad: $\bigstar\bigstar\bigstar$.
\end{prioritario}

\textbf{Paginas}: 60 ~|~ \textbf{Audiencia}: RF / comunicaciones.

\section{Contenido}

\begin{enumerate}[leftmargin=2em]
  \item \textbf{Bandas usadas en CubeSat}: VHF (143 MHz), UHF (433-460 MHz),
        S-Band (2.2-2.4 GHz), X-Band (8 GHz), Ka-Band (26 GHz).
  \item \textbf{Link budget}: potencia tx, ganancia antenas, free space
        loss, ruido, margen.
  \item \textbf{Modulacion}: BPSK, QPSK, GMSK, FSK, OFDM.
  \item \textbf{Codificacion}: FEC, Reed-Solomon, Turbo codes, LDPC,
        convolutional.
  \item \textbf{Protocolos}: AX.25, CCSDS Space Packet, KISS, CSP (CubeSat
        Space Protocol).
  \item \textbf{Antenas en CubeSat}: dipolo (UHF), patch (S-Band), helical,
        omni vs directional.
  \item \textbf{Despliegue de antenas}: burn-wire deployment, dipole
        unfold.
  \item \textbf{Ground station}: receptor, decoder, doppler tracking,
        rotator azimuth/elevation, SatNOGS.
  \item \textbf{Coordinacion con IARU}: registracion de frecuencia,
        coordinacion internacional, licencias amateur radio.
  \item \textbf{Componentes comerciales}: ISIS UHF/VHF, EnduroSat S-Band,
        Tethers Unlimited X-Band, GomSpace.
\end{enumerate}

% =====================================================================
\chapter{Manual M5: Flight software y OBC}

\begin{prioritario}
Prioridad: $\bigstar\bigstar\bigstar$.
\end{prioritario}

\textbf{Paginas}: 60 ~|~ \textbf{Audiencia}: software embebido.

\section{Contenido}

\begin{enumerate}[leftmargin=2em]
  \item \textbf{Arquitecturas de OBC}: microcontrolador vs Linux SoC,
        cuando usar cada uno.
  \item \textbf{Sistemas operativos}: FreeRTOS, RTEMS, Zephyr, Linux,
        VxWorks (no en CubeSat usualmente).
  \item \textbf{Frameworks de flight software}: NASA cFS, F' (FPrime), KubOS,
        BRASS.
  \item \textbf{Tareas vs interrupciones}: prioridades, scheduling.
  \item \textbf{Comunicacion inter-tareas}: queues, semaphores, pub/sub.
  \item \textbf{Telemetria y comandos}: CCSDS, beacon, housekeeping.
  \item \textbf{FDIR}: deteccion, aislamiento, recuperacion de fallos.
  \item \textbf{Watchdog hardware y software}.
  \item \textbf{OTA updates}: como actualizar firmware en orbita con rollback.
  \item \textbf{Versiones y configuration management}: git, branches,
        releases.
  \item \textbf{Casos especificos}: el daemon del INTISAT, CCSDS sobre I2C,
        chunking.
\end{enumerate}

% =====================================================================
\chapter{Manual M6: Control termico}

\textbf{Paginas}: 40 ~|~ \textbf{Audiencia}: termico / materiales.
~|~ Prioridad: $\bigstar\bigstar$.

\section{Contenido}

\begin{enumerate}[leftmargin=2em]
  \item \textbf{Modos de transferencia de calor}: conduccion, radiacion
        (no conveccion en vacio).
  \item \textbf{Calculo de balance termico}: $Q_{\text{in}} = Q_{\text{out}}$
        en equilibrio.
  \item \textbf{Propiedades de superficie}: emisividad, absorbtividad,
        ratio $\alpha/\epsilon$.
  \item \textbf{Materiales}: pinturas, MLI (Multi-Layer Insulation),
        thermal tape, radiators.
  \item \textbf{Control activo vs pasivo}: heaters, conductive paths.
  \item \textbf{Componentes sensibles}: bateria (debe estar $> 0$ °C para
        cargar), sensor de imagen (drift termico).
  \item \textbf{Analisis termico}: Thermal Desktop, FreeCAD-FEM,
        NX Thermal, ESATAN.
  \item \textbf{TVAC test}: ciclos en vacio, gradientes, soak.
  \item \textbf{Disipacion en el Pi 5}: pad termico al frame, limite de
        frecuencia.
\end{enumerate}

% =====================================================================
\chapter{Manual M7: Integracion y test campaign}

\begin{prioritario}
Prioridad: $\bigstar\bigstar\bigstar$.
\end{prioritario}

\textbf{Paginas}: 60 ~|~ \textbf{Audiencia}: AIT (Assembly, Integration,
Test).

\section{Contenido}

\begin{enumerate}[leftmargin=2em]
  \item \textbf{Filosofia de testing en aeroespacial}: test as you fly,
        protoflight, qualification model, flight model.
  \item \textbf{Estandares}: NASA GEVS, ECSS-E-ST-10-03, MIL-STD-810.
  \item \textbf{Test funcional}: smoke test, comm test, end-to-end.
  \item \textbf{Test de vibracion}: sine sweep, random, shock,
        niveles (3 g RMS tipico CubeSat).
  \item \textbf{Test termico vacuum (TVAC)}: bake-out, ciclos, soak,
        funcional caliente y frio.
  \item \textbf{Test de vacio}: outgassing, vacuum chamber alone.
  \item \textbf{Test EMC/EMI}: emisiones radiadas, conducidas,
        susceptibilidad.
  \item \textbf{Test de radiacion}: TID con Co-60, SEU con protones, TLD.
  \item \textbf{Test funcional final}: 100 scans, 168 h burn-in.
  \item \textbf{Mass properties}: balance, CG, inercia.
  \item \textbf{Cleanliness}: clean room class 100K.
  \item \textbf{Reviews}: PDR, CDR, TRR, PSR, MRR.
\end{enumerate}

% =====================================================================
\chapter{Manual M8: Operaciones y ground station}

\textbf{Paginas}: 50 ~|~ \textbf{Audiencia}: operadores de mision.
~|~ Prioridad: $\bigstar\bigstar$.

\section{Contenido}

\begin{enumerate}[leftmargin=2em]
  \item \textbf{Ciclo de operaciones diario}: passes, comandos, downlink.
  \item \textbf{Diseño de ground station}: antena, rotator, SDR,
        software (gpredict, satnogs).
  \item \textbf{Planificacion de pases}: STK, GMAT, GPredict.
  \item \textbf{Mission control software}: OpenMCT, COSMOS, Yamcs.
  \item \textbf{Telemetria}: parsing CCSDS, dashboards, alarmas.
  \item \textbf{Comandos uplink}: secuencias, ACK, retries.
  \item \textbf{Anomaly response}: procedures, escalacion.
  \item \textbf{Flight rules}: que se puede y no se puede hacer.
  \item \textbf{Decomissioning}: deorbit, passivacion, end-of-life.
  \item \textbf{Coordinacion con autoridades}: ANATEL, IARU, FCC.
\end{enumerate}

% =====================================================================
\chapter{Manual M9: Project management y system engineering}

\textbf{Paginas}: 40 ~|~ \textbf{Audiencia}: lider de proyecto.
~|~ Prioridad: $\bigstar$ (mas tarde).

\section{Contenido}

\begin{enumerate}[leftmargin=2em]
  \item \textbf{Fases NASA}: Pre-A, A, B, C, D, E, F.
  \item \textbf{Reviews tecnicas}: SRR, PDR, CDR, TRR, FRR, PSR, ORR.
  \item \textbf{Requirements engineering}: como redactar y rastrear.
  \item \textbf{V-model de verificacion y validacion}.
  \item \textbf{Configuration management}: baseline, change requests.
  \item \textbf{Risk management}: matriz probabilidad/impacto.
  \item \textbf{Cost estimation}: Cost models para CubeSat.
  \item \textbf{Schedule}: Gantt, dependencias, critical path.
  \item \textbf{Team organization}: roles, communicacion, herramientas
        (Jira, Notion, Trello).
\end{enumerate}

% =====================================================================
\chapter{Manual M10: CM5 y carrier para INTISAT}

\begin{prioritario}
Prioridad: $\bigstar\bigstar\bigstar$ --- es el siguiente paso concreto
para el INTISAT.
\end{prioritario}

\textbf{Paginas}: 50 ~|~ \textbf{Audiencia}: hardware del payload.

\section{Contenido}

\begin{enumerate}[leftmargin=2em]
  \item \textbf{Diferencias Pi 5 vs CM5}.
  \item \textbf{Conector mezzanine}: DF40 200-pin, voltajes, signal map.
  \item \textbf{Diseño del carrier}: KiCad, layout, layers.
  \item \textbf{Componentes en el carrier}: LDO, fusible, INA219,
        TXS0108E, conectores.
  \item \textbf{Compatibilidad PC-104}: pinout, footprint.
  \item \textbf{Power-on sequence}.
  \item \textbf{Flash de eMMC}: rpiboot, imagen pre-configurada.
  \item \textbf{Conectores externos}: CSI-2 para OV5647, SPI para OLED.
  \item \textbf{Test de banco}: bring-up procedure.
  \item \textbf{Validacion ambiental}: vibracion y TVAC del carrier.
\end{enumerate}

% =====================================================================
\chapter{Orden de lectura sugerido}

\section{Para un miembro nuevo (full onboarding)}

\textbf{Semana 1}: M0 (CubeSat 101) + lo basico del proyecto en
\texttt{contexto\_sesion/}.

\textbf{Semana 2}: Manual\_Electronica\_Practica.pdf (ya existe) + M2 (EPS)
fundamentos.

\textbf{Semana 3}: Raspberry\_Pi\_5\_Profundo.pdf (ya existe) + M5
(Flight Software).

\textbf{Semana 4}: M4 (Comms) + M7 (Testing) capitulos 1-3.

Despues, especializacion segun rol.

\section{Para un ingeniero electronico de payload}

M0 $\to$ Manual\_Electronica\_Practica $\to$ Raspberry\_Pi\_5\_Profundo $\to$
M2 (EPS) $\to$ M10 (CM5 carrier) $\to$ M7 (Testing).

\section{Para un operador de mision}

M0 $\to$ M4 (Comms) $\to$ M8 (Ops) $\to$ Referencias\_OpenSource\_CubeSat
(OpenMCT) $\to$ Bibliografia\_AI\_Espacial.

\section{Para un lider de proyecto}

M0 $\to$ M9 (Project management) $\to$ M7 (Testing) $\to$ resto en panorama.

% =====================================================================
\chapter{Cronograma de produccion sugerido}

\section{Si el equipo es de 4 personas, 15 semanas}

\begin{center}
\begin{tabular}{lll}
\toprule
\textbf{Sprint} & \textbf{Manual a escribir} & \textbf{Responsable} \\
\midrule
1--2  & M0 CubeSat 101                            & lider tecnico \\
3--4  & M2 EPS (paralelo con M5)                  & elec/potencia \\
3--4  & M5 Flight software (paralelo con M2)      & software \\
5--6  & M4 Comms (paralelo con M10)               & RF \\
5--6  & M10 CM5 carrier (paralelo con M4)         & hw payload \\
7--8  & M7 Testing                                 & AIT \\
9--10 & M1 Mecanica y M3 ADCS (paralelo)          & mec y control \\
11--12 & M6 Termico y M8 Operaciones (paralelo)    & termico y ops \\
13--14 & M9 Project management                     & lider proyecto \\
15    & Revision cruzada, integracion              & todos \\
\bottomrule
\end{tabular}
\end{center}

\section{Si son menos personas o menos tiempo}

\textbf{Si solo hay 2 personas}: enfocarse en los $\bigstar\bigstar\bigstar$
unicamente. Total: M0 + M2 + M4 + M5 + M7 + M10 = 6 manuales, $\sim 340$ pag,
$\sim 9$ semanas.

\textbf{Si solo hay 1 persona escribiendo}: $\sim 30$ semanas para los 6
prioritarios. Mejor reclutar mas.

% =====================================================================
\chapter{Convenciones de estilo}

Para mantener consistencia entre manuales escritos por distintas personas:

\begin{enumerate}[leftmargin=2em]
  \item \textbf{LaTeX class}: \texttt{book} para los grandes ($> 50$ pag),
        \texttt{report} para los medianos.
  \item \textbf{Idioma}: español (es-noshorthands, sin acentos en codigo).
  \item \textbf{Cajas de color} consistentes:
    \begin{itemize}[leftmargin=1em]
      \item Verde: \texttt{intuibox} (analogias, intuiciones)
      \item Gris: \texttt{calculo} (formulas, derivaciones)
      \item Azul: \texttt{payload} (conexion con INTISAT)
      \item Roja: \texttt{cuidado} (errores comunes)
      \item Naranja: \texttt{ejercicio} (practica)
    \end{itemize}
  \item \textbf{Cada capitulo}: la misma plantilla (por que importa,
        intuicion, mate, conexion, cuidado, ejercicios).
  \item \textbf{Bibliografia}: minimo 5 referencias por manual.
  \item \textbf{Apendices}: tablas de consulta + roadmap de aprendizaje.
\end{enumerate}

\begin{nota}
\textbf{Plantilla base}: tomar \texttt{Manual\_Electronica\_Practica.tex}
como template y reemplazar el contenido. Asi todos los manuales tendran
el mismo \texttt{look \& feel}.
\end{nota}

% =====================================================================
\chapter{Estimacion de esfuerzo}

\section{Promedio por pagina}

Para un manual tecnico de calidad (escrito desde cero con investigacion):

\begin{itemize}[leftmargin=2em]
  \item Investigacion previa: 30 min/pagina.
  \item Redaccion: 60 min/pagina.
  \item Diagramas / tablas: 30 min/pagina.
  \item Revision y correccion: 20 min/pagina.
  \item \textbf{Total}: $\sim 2.3$ h/pagina.
\end{itemize}

\section{Aplicado a la coleccion}

\begin{center}
\begin{tabular}{lcc}
\toprule
\textbf{Manual} & \textbf{Pag} & \textbf{Horas} \\
\midrule
M0 CubeSat 101                       & 50 & 115 \\
M1 Mecanica                          & 50 & 115 \\
M2 EPS                                & 60 & 138 \\
M3 ADCS                               & 60 & 138 \\
M4 Comms                              & 60 & 138 \\
M5 Flight software                    & 60 & 138 \\
M6 Termico                            & 40 & 92  \\
M7 Testing                            & 60 & 138 \\
M8 Operaciones                        & 50 & 115 \\
M9 PM/SE                              & 40 & 92  \\
M10 CM5 carrier                       & 50 & 115 \\
\midrule
\textbf{Total}                       & \textbf{580} & \textbf{1334 h} \\
\bottomrule
\end{tabular}
\end{center}

$\sim 1334$ h $= 167$ dias-persona de 8 h $= 33$ semanas-persona.
Con un equipo de 4 personas: $\sim 8$ semanas reales (dedicacion plena).
Con dedicacion parcial (2 h/dia): $\sim 6$ meses.

% =====================================================================
\chapter{Plan de arranque concreto}

\section{Primer manual a escribir: M0 CubeSat 101}

\textbf{Justificacion}: es el que mas valor genera (todo el equipo se
beneficia), no requiere prerequisitos, y es el que mejor consolida el
vocabulario tecnico que todos deberian compartir.

\textbf{Tiempo estimado}: 115 h, 1 persona dedicada full-time durante 3
semanas, o equipo en sprint.

\section{Segundo manual: M5 (Flight Software) o M2 (EPS)}

Segun cual subsistema esta mas adelantado en INTISAT, priorizar:

\begin{itemize}[leftmargin=2em]
  \item Si el payload de microscopia esta avanzado (lo que vemos) $\to$ M5.
  \item Si el sistema de potencia esta en bring-up $\to$ M2.
\end{itemize}

\section{Tercer manual: M10 CM5 carrier}

Justifica una vez que se decide migrar el payload a CM5. Es muy especifico
pero da el siguiente nivel concreto al INTISAT.

% =====================================================================
\chapter{Conclusion}

\begin{tip}
\textbf{Decision sugerida}: empezar por \textbf{M0 CubeSat 101} esta
semana. Es de alto impacto, no bloquea otras tareas, y genera material
inmediatamente util para onboarding del equipo. Asignar a una persona
dedicada o repartir capitulos entre el equipo (cada uno escribe 1
capitulo de 5 paginas en 1 semana $=$ 50 paginas en 1 semana del equipo
de 10).
\end{tip}

\textbf{Una vez completo M0}, el equipo decide segun prioridades del
proyecto:
\begin{itemize}[leftmargin=2em]
  \item Si la prioridad es \textbf{integrar el payload al INTISAT}: M2, M5, M10.
  \item Si la prioridad es \textbf{prepararse para vuelo}: M7, M8.
  \item Si la prioridad es \textbf{onboarding} de nuevos miembros: solo M0
        alcanza.
\end{itemize}

Esta coleccion no se hace en un mes. Es un activo de mediano plazo
(6 meses con dedicacion parcial). Pero \textbf{cada manual completado vale
por si mismo}: es util desde el dia que se publica, no espera a estar
toda la serie completa.

\end{document}
